\fchapter{Ej 6 Pág 77}

\fsection{\boldit{\textcolor{waveBlue2}{Actividad.}}} 
\subsection*{El sector hotelero es uno de los más comprometidos con el uso eficiente de la energía y la búsqueda de soluciones sostenibles. Enumera algunas medidas que se pueden llevar a cabo en cualquier establecimiento, tanto en las habitaciones como en las zonas comunes, que ayuden a la consecución del ODS 7.}

\begin{solution}
	\begin{enumerate}	
	    \item \textbf{En habitaciones:}
		\begin{itemize}

		    \item \textbf{Instalar llaves maestras que activen la luz y climatización solo con la tarjeta de la habitación.}

		    \item \textbf{Usar iluminación LED y sensores de presencia en baños y pasillos.}

		    \item \textbf{Implementar termostatos inteligentes para un control eficiente de la calefacción y aire acondicionado.}

		\end{itemize}

	    \item \textbf{En zonas comunes y operaciones:}
		\begin{itemize}

		    \item \textbf{Optimizar la climatización de zonas como el lobby, gimnasio o pasillos.}

		    \item \textbf{Instalar paneles solares para generar energía renovable.}

		    \item \textbf{Utilizar sistemas de gestión energética centralizada para monitorizar y controlar el consumo en tiempo real.}

		\end{itemize}

	    \item \textbf{En equipamiento y mantenimiento:}

		\begin{itemize}
		    \item \textbf{Adquirir electrodomésticos y equipos (lavadoras, cocinas) con certificación de alta eficiencia energética (A+++).}
		    \item \textbf{Establecer un programa de mantenimiento preventivo de calderas y sistemas de climatización.}
		    \item \textbf{Cambiar a vehículos eléctricos o híbridos para los servicios del hotel (transfers, mantenimiento).}
		\end{itemize}

	\end{enumerate}






\end{solution}
